\documentclass[11pt]{article}
\usepackage[margin=2.4cm]{geometry}
\usepackage{amsmath,amssymb,amsthm,booktabs}
\usepackage[T1]{fontenc}
\usepackage{hyperref}
\sloppy
\newtheorem{theorem}{Theorem}
\newtheorem{lemma}{Lemma}
\newtheorem{conjecture}{Conjecture}
\newtheorem{question}{Question}
\theoremstyle{definition}
\newtheorem{remark}{Remark}
\title{An inequality between determinant and signature for $L$-space knots\\[2mm]
\large Machine-generated conjecture, its checks, a partial proof, and the open case}
\date{2026-09-19}
\begin{document}
\maketitle

\begin{abstract}
We consider the inequality $\det(K) \le 1 + |\sigma(K)|$ for $L$-space knots. It was produced by a program that fits
linear inequalities to the KnotInfo database and discards what its automatic checks refute. It survived those checks
and the computations of Section 3. We prove it for every $L$-space knot that is an iterated torus knot, in
particular for all algebraic knots. The hyperbolic case is open, and it is where the inequality is sharp:
equality holds on every member tested of an infinite family of hyperbolic $L$-space knots, the first of which is not
braid positive. Among the 632 census $L$-space knots it holds only on that first member.
Sections 2 and 3 give the provenance and the computations, so that the proof in Sections 5 and 6 can be checked
without our files; the one lemma we prove from scratch rather than cite is additionally formalised in Lean 4
with no \texttt{sorry}.
\end{abstract}

\section{Statement and conventions}

Signs follow KnotInfo: the right-handed trefoil has $\sigma = -2$, $s = 2$, $\tau = 1$. A knot $K \subset S^3$ is an
\emph{$L$-space knot} if some positive Dehn surgery on $K$ yields an $L$-space. Positive torus knots then have
$\sigma < 0$. We do not assume a sign for a general $L$-space knot: $\sigma \le 0$ is carried through the induction of
Theorem~\ref{thm} for the knots treated there, and it held on every $L$-space knot we computed (all 632 census knots
and all 68 knots of the families of Section~3). Write $\det(K) = |\Delta_K(-1)|$ and $g = g(K)$ for the Seifert genus.

\begin{conjecture}[K33]\label{conj}
For every $L$-space knot $K$, \[ \det(K) \;\le\; 1 + |\sigma(K)| . \]
\end{conjecture}

We use the following standard facts. An $L$-space knot is fibered, strongly quasipositive, and satisfies
$g = g_4 = \tau$ (Ozsv\'ath--Szab\'o, Ni, Hedden); its Alexander polynomial has the form
$\Delta_K(t) = \sum_{k=0}^{2n} (-1)^k t^{a_k}$ with $a_0 = g > a_1 > \dots > a_{2n} = -g$ (Ozsv\'ath--Szab\'o).
Consequently $|\sigma(K)| \le 2g$ and $\det(K) \le 2g+1$.

\section{Provenance}

The inequality was not conjectured by a person. It is candidate K33 of 42 statements emitted by \texttt{ragusa}, an
engine that enumerates bounds of the form (property) $\Rightarrow$ (target $\le c \cdot$ feature $+ k$) over a table of
invariants, here KnotInfo \cite{KnotInfo} restricted to prime knots with 3 to 13 crossings, and reports what survives its filters
together with the evidence behind each survivor. K33 was found on the \texttt{l\_space} column, where the table
contains 11 knots; those knots were the whole of its initial evidence.

The other 41 candidates were triaged in the same session: 18 were refuted by explicit knots, 13 turned out to be known
theorems or true by definition, 3 had one clause refuted and one left open, and 6 remain open but rest on weak
evidence. Two remained open with evidence beyond the table: K33, and $u = g$ for $L$-space knots.

\section{Verification history}

All computations use SnapPy (knot Floer homology, for the $L$-space property and $\tau$), Sage (Seifert matrices,
signatures, determinants) and exact formula implementations; scripts are listed in Section~\ref{sec:repro}.

\begin{center}\small
\begin{tabular}{p{5.1cm}rrp{5.4cm}}
\toprule
family of $L$-space knots & tested & viol. & how the invariants are known \\
\midrule
iterated torus knots & 11728 & 0 & exact: cabled Alexander polynomial, Litherland's formula \\
SnapPy census $L$-space knots (all) & 632 & 0 & Seifert matrix of a positive braid word \\
census knots, independent route & 271 & 0 & knot Floer homology, Sage signature \\
Baker--Kegel, Himeno, $(2,q)$-cables & 68 & 0 & Seifert matrix of the braid word \\
one-bridge braids $B(w,b,t)$, $w \le 15$ (parameter triples) & 274 & 0 & Seifert matrix of the braid word \\
twisted torus knots $K(p,q;r,s)$ (distinct knots) & 87 & 0 & knot Floer homology, then exact Seifert signature \\
random braid words (distinct knots) & 46 & 0 & knot Floer homology, then exact Seifert signature \\
\bottomrule
\end{tabular}
\end{center}

\noindent The census braid words are those published by Baker and Kegel (arXiv:2203.12013, Appendix), which cover all
632 census $L$-space knots. Every closure was checked to be the named census knot: 631 of 632 by SnapPy's isometry test
directly, the remaining one after retriangulation and at higher precision. All 632 signatures are exact, computed by
congruence diagonalisation over $\mathbb{Q}$ from the integral matrix $V + V^{\mathsf T}$; they agree with an earlier
floating-point computation on all 632 knots and with Sage's own exact signature on the 511 knots where that finished.
The one-bridge braids are $L$-space knots by \cite{GLV}; for the 40 smallest this was also confirmed directly by knot
Floer homology, with the genus matching the braid-surface genus.

\paragraph{Generated families, counted by knot.} Three rows of the table come from parameters or words rather than
from a list: the one-bridge braids, the twisted torus knots $K(p,q;r,s)$, the closures of
$(\sigma_1 \cdots \sigma_{p-1})^q (\sigma_1 \cdots \sigma_{r-1})^s$, and the closures of random braid words. Different
parameters often close to the same knot, and an earlier version of this note counted diagrams, which overstated these
families several times over. Counted by knot type, exactly for hyperbolic knots through the isometry signature of the
complement and approximately otherwise, the 274 one-bridge parameter triples give 195 knots, 149 of them
hyperbolic and at least 37 of those outside the SnapPy census; the 549 twisted torus records give 87
knots, 42 hyperbolic; and the 14452 random records give 46 knots, 17 hyperbolic. Every hyperbolic
knot in the last two families is a census knot, so they confirm the census rows rather than extend them; the hyperbolic
evidence beyond the census is the one-bridge braids and the Baker--Kegel and Himeno families. In the twisted torus and
random families the $L$-space property was decided by each knot's own knot Floer homology, and the determinant was
computed twice, from the knot Floer gradings and from the Seifert matrix, agreeing on all 15001 records.

\paragraph{Independent check.} A language model with no access to our computations re-derived the lattice count and
recomputed, independently: the six torus signatures used here, Lemma~\ref{lem:L} and the bound on $N_<$ for all coprime
$p \le 24$, $q < 120$, the closed forms and their reductions, the list of leftover small pairs in the induction, the
induction step under worst-case hypotheses, and the identity \eqref{eq:det} on three examples. It found no errors. The
gaps it flagged (an unsupported sign assertion, boundary points in the count, the $p<q$ convention, the nontriviality
hypothesis in (F4), the derivation of the symmetry of $S$, missing references) are addressed in this version. A third model corrected the proof of Lemma~\ref{lem:L}, where an earlier version claimed an
equivalence that is only a sufficient condition, and asked for the hypotheses and chirality conventions now stated;
the identity $\sigma = 4N_{<} - 2g$ was verified on all coprime $p \le 15$, $q < 50$, where $N_{<}/g \le 0.234$, so
the branch proved here is the operative one. A second
model, asked to attack the open case, contributed the congruence of Section~5, the one-bridge braid family of the
table, and the rigidity observation of Section~7; its computations of the one-bridge braids (274 parameter triples, 0 violations,
70 equalities) and of the Baker--Kegel invariants were reproduced independently here.

\paragraph{The hypothesis is not redundant.} Evaluating the same inequality on census knots whose braid word is
positive or negative but which are \emph{not} $L$-space knots produces 25 violations. So the $L$-space condition does
the work; this is unlike several of the 42 candidates, whose hypotheses turned out to be irrelevant.

\paragraph{Sharpness.} Equality $\det = 1 + |\sigma|$ holds for $T(2,n)$, for the $(2,q)$-cables of the trefoil, and,
among the families tested beyond the census, for 62 of 68 knots, including every Baker--Kegel $K_n$ with
$n = 1, \dots, 25$. These are hyperbolic: $K_1$ is the census knot o9\_30634, the one census $L$-space knot that is
not braid positive, and $K_n$ for $n \ge 2$ is not known to be braid positive. Of the 632 census $L$-space knots,
exactly 1 attains equality, and it is $K_1$.

\section{A reformulation}

For an $L$-space knot let $S \subset \mathbb{Z}_{\ge 0}$ be its staircase set (the \emph{formal semigroup}; for an
algebraic knot it is the semigroup of the singularity, but for a general $L$-space knot it need not be closed under
addition), the complement of the $g$ gaps, determined by $t^{g}\Delta_K(t) = (1-t)\sum_{s \in S,\, s < 2g} t^{s} + t^{2g}$. Evaluating at $t = -1$ and using the symmetry
$s \leftrightarrow 2g-1-s$ between $S \cap [0,2g)$ and the gaps, which follows from $\Delta_K(t) = \Delta_K(t^{-1})$
and gives $|S \cap [0,2g)| = g$ as well, one obtains

\begin{equation}\label{eq:det}
 \det(K) = |2D+1|, \qquad D = \#\{\text{odd gaps}\} - \#\{\text{even gaps}\}.
\end{equation}

Hence Conjecture~\ref{conj} is equivalent to $|\sigma| \ge 2D$ when $D \ge 0$, and to $|\sigma| \ge 2|D| - 2$ when
$D < 0$. Identity~\eqref{eq:det} was checked on 11728 iterated torus knots with 0 mismatches.

\begin{remark}
The cruder route fails. Since the coefficients of $\Delta_K$ are $\pm 1$, $\det(K)$ is at most the number of nonzero
coefficients, so Conjecture~\ref{conj} would follow from $|\sigma| \ge \#\{\text{terms}\} - 1$. That inequality is false:
it fails on 401 of the 11728 iterated torus knots tested, and the failures are exactly the knots where
Conjecture~\ref{conj} is sharp. Any proof must therefore see the cancellation in $\Delta_K(-1)$.
\end{remark}

\section{An arithmetic refinement}

For any knot, $\Delta_K(-1) \equiv 1 \pmod 4$ and $\operatorname{sign}\Delta_K(-1) = (-1)^{\sigma/2}$ (Murasugi),
so $\det(K) \equiv (-1)^{\sigma/2} \pmod 4$ and therefore

\begin{equation}\label{eq:cong}
 1 + |\sigma(K)| - \det(K) \equiv 0 \pmod 4 .
\end{equation}

Checked on all 1523 $L$-space knots computed here with an exact signature and determinant: 0 failures. Two consequences. First, a counterexample to Conjecture~\ref{conj} must satisfy
$\det(K) \ge |\sigma(K)| + 5$; the inequality cannot fail narrowly. Second, combining \eqref{eq:cong} with
\eqref{eq:det}, the \emph{signature defect} $2g - |\sigma|$ of an $L$-space knot satisfies
$2g - |\sigma| \equiv 0 \pmod 4$ when $D \ge 0$, and Conjecture~\ref{conj} becomes

\[ 2g - |\sigma| \;\le\; 4\,\#\{\text{even gaps}\} \qquad (D \ge 0), \]

both sides being multiples of $4$. (For $D < 0$ the corresponding statement is
$2g - |\sigma| \le 4\,\#\{\text{odd gaps}\} + 2$.) In particular the conjecture holds automatically for every
$L$-space knot with $\det \le 3$ and for every definite one.

\section{Proof for iterated torus knots}

Throughout, $p \ge 2$ and $q$ are coprime, $T(p,q)$ is the positive torus knot with $g(T(p,q)) = (p-1)(q-1)/2$, and
$K_{p,q}$ denotes the $(p,q)$-cable of $K$ (winding number $p$). We use:

\begin{itemize}
\item[(F1)] $\Delta_{K_{p,q}}(t) = \Delta_K(t^p)\,\Delta_{T(p,q)}(t)$; since $\Delta_K(1) = 1$ this gives
  $\det(K_{p,q}) = \det(K)\det(T(p,q))$ for $p$ odd and $\det(K_{p,q}) = \det(T(p,q))$ for $p$ even.
\item[(F2)] Litherland's cabling formula $\sigma(K_{p,q}) = \sigma_{(-1)^p}(K) + \sigma(T(p,q))$, so
  $\sigma(K_{p,q}) = \sigma(K) + \sigma(T(p,q))$ for $p$ odd and $= \sigma(T(p,q))$ for $p$ even.
\item[(F3)] $\det(T(p,q)) = 1$ if $p,q$ are both odd, $= p$ if $q$ is even, $= q$ if $p$ is even.
\item[(F4)] (Hedden; Hom) for $K$ nontrivial, $K_{p,q}$ is an $L$-space knot if and only if $K$ is one and
  $q \ge p(2g(K)-1)$.
\end{itemize}

(F1)--(F3) were verified on 4088 cable instances, with 0, 0 and 0 mismatches respectively.

\begin{lemma}\label{lem:L}
For integers $2 \le p < q$ with $\gcd(p,q) = 1$, \quad
$\sigma(T(p,q)) \le -g(T(p,q))$, where $g(T(p,q)) = \tfrac{(p-1)(q-1)}{2}$; in particular $|\sigma| \ge g$.
\end{lemma}

\begin{proof}
We use the lattice count for torus knot signatures (Brieskorn; Gordon--Litherland--Murasugi \cite{GLM};
Litherland \cite{Lith}), in the normalisation $\sigma(T(2,3)) = -2$ of Section 1:
$\sigma(T(p,q)) = N_{\mathrm{out}} - N_{\mathrm{in}}$ over the
grid $G = \{1,\dots,p-1\}\times\{1,\dots,q-1\}$, where $(i,j)$ is counted in $N_{\mathrm{in}}$ when
$\tfrac12 < \tfrac ip + \tfrac jq < \tfrac32$. The involution $(i,j) \mapsto (p-i,q-j)$ of $G$ exchanges the two
outside regions, so $N_{\mathrm{out}} = 2N_{<}$ with $N_{<} = \#\{(i,j) \in G : \tfrac ip + \tfrac jq < \tfrac12\}$.
No lattice point of $G$ lies on a boundary: $\tfrac ip + \tfrac jq \in \{\tfrac12, \tfrac32\}$ gives
$2(iq + jp) \in \{pq, 3pq\}$, so $p \mid 2i$ and $q \mid 2j$, forcing $2i = p$ and $2j = q$ with $p$ and $q$ both even,
contradicting $\gcd(p,q) = 1$. As $N_{\mathrm{in}} + N_{\mathrm{out}} = 2g$, the count reads $\sigma(T(p,q)) = 4N_{<} - 2g$, so $|\sigma| \ge g$
holds if $N_{<} \le g/4$, and also if $N_{<} \ge 3g/4$. We prove the first branch,
$N_{<} \le (p-1)(q-1)/8$, which gives $\sigma(T(p,q)) \le -g < 0$.

For fixed $i$ the number of admissible $j$ is $\lceil f(i)\rceil - 1 \le f(i)$ with $f(i) = \tfrac q2 - \tfrac{iq}{p}$,
and $f(i) > 0$ exactly for $i < p/2$. Summing over those $i$:
if $p$ is even, with $m = p/2-1$, one gets $N_{<} \le q(p-2)/8$, and $q(p-2) \le (p-1)(q-1)$ reduces to $q \ge p-1$;
if $p$ is odd, with $m = (p-1)/2$, one gets $N_{<} \le q(p-1)^2/(8p)$, and $q(p-1)/p \le q-1$ reduces to $p \le q$.
Both hold because $q > p$.
\end{proof}

\paragraph{The lattice inequality is machine-checked.} The combinatorial core of this proof, the bound
$N_{<} \le (p-1)(q-1)/8$ for all integers $2 \le p < q$, has been formalised in Lean 4 against Mathlib
and carries no \texttt{sorry}: \texttt{\#print axioms} reports only \texttt{propext},
\texttt{Classical.choice} and \texttt{Quot.sound}. Coprimality is not needed for it. The formalisation
covers the column decomposition of the grid count, the per-column bound, the cut at $(p-1)/2$ in both
parities, the Gauss sum, the closed forms, and the two final comparisons; it does not cover the topology
the lemma is used with, namely the lattice count for $\sigma(T(p,q))$, which is cited. So the part of
Lemma~\ref{lem:L} proved here from scratch, rather than cited, is the part that is now checked by a
machine. The file is \texttt{K33Lattice/Basic.lean}.

\noindent Every step of this proof was also checked numerically on 429 pairs with $p \le 15$, $q < 60$ (0 failures), and the conclusion
holds on all 2166 pairs with $p \le 20$, $q < 200$. Lemma~\ref{lem:L} says $-\sigma \ge b_1/2$ for torus knots, which is stronger than the linear bounds known for
arbitrary positive braid links (Feller proved a linear bound; Greene--Liechti give $-\sigma \ge b_1/4$, see the
abstract of \cite{GL}). We did not find Lemma~\ref{lem:L} in the literature; it may be folklore.

\begin{theorem}\label{thm}
Every $L$-space knot which is an iterated torus knot satisfies $\det(K) \le 1 + |\sigma(K)|$. In particular this holds
for all positive torus knots, hence for their mirrors since $\det$ and $|\sigma|$ are mirror invariant, and for all
algebraic knots, which are iterated torus knots with positive parameters \cite{EN} and hence $L$-space knots.
\end{theorem}

\begin{proof}
The induction is on the cabling depth, with torus knots at depth one; by (F4) applied to the companion (Hom's
``only if''), every companion of an $L$-space iterated torus knot is again an $L$-space knot, so the hypothesis is
available at each step. We prove simultaneously that $\det(K) \le 1 + |\sigma(K)|$ and $\sigma(K) < 0$. For torus knots
we may assume $p < q$, since $T(p,q) = T(q,p)$; in a cable $K_{p,q}$ of a nontrivial $L$-space knot, (F4) gives
$q \ge p(2g_K - 1) \ge p$, and $\gcd(p,q) = 1$ excludes $q = p$.

\emph{Base case.} For $T(p,q)$, $\sigma < 0$ is part of Lemma~\ref{lem:L}, and: if $p,q$ are both odd then $\det = 1$. If $q$ is even then $\det = p$ and
Lemma~\ref{lem:L} gives $|\sigma| \ge 3(p-1)/2$ (using $q \ge 4$), whence $1 + |\sigma| \ge p$. If $p$ is even then
$\det = q$; for $p = 2$ we have $|\sigma(T(2,q))| = q-1$, so $1 + |\sigma| = q$, with equality; for $p \ge 4$,
Lemma~\ref{lem:L} gives $|\sigma| \ge 3(q-1)/2 \ge q-1$.

\emph{Induction step.} Let $K$ be an $L$-space knot with $\det(K) \le 1 + A$, where $A = |\sigma(K)| \le 2g_K$, and let
$K_{p,q}$ be an $L$-space cable, so $q \ge p(2g_K-1)$ by (F4). Put $B = |\sigma(T(p,q))|$. By the induction hypothesis $\sigma(K) < 0$, and
$\sigma(T(p,q)) < 0$ by Lemma~\ref{lem:L}, so (F2) gives $\sigma(K_{p,q}) < 0$ and $|\sigma(K_{p,q})| = A + B$ for $p$
odd, $= B$ for $p$ even; this also propagates the sign claim.

If $p$ is even, then $\det(K_{p,q}) = \det(T(p,q)) = q$ by (F1) and (F3), and $q \le 1 + B$ was shown in the base case.

If $p$ and $q$ are both odd, then $\det(K_{p,q}) = \det(K) \le 1 + A \le 1 + A + B$.

If $p$ is odd and $q$ even, then $\det(K_{p,q}) = p\det(K) \le p(1+A)$, and the claim $p(1+A) \le 1 + A + B$ is
$(p-1)(1+A) \le B$. By Lemma~\ref{lem:L} it suffices that $1 + A \le (q-1)/2$, i.e. $q \ge 2A+3$, and since
$A \le 2g_K$ it suffices that $q \ge 4g_K+3$. As $p \ge 3$, (F4) gives $q \ge 6g_K-3$, which is $\ge 4g_K+3$ exactly
when $g_K \ge 3$. For $g_K \in \{1,2\}$ the uncovered pairs are finite: for $g_K = 1$ they are $(p,q) = (3,4)$ and
$(5,6)$, where $B = 6 \ge 6$ and $B = 16 \ge 12$; for $g_K = 2$ it is $(3,10)$, where $B = 14 \ge 10$.
\end{proof}

\noindent The induction step was also verified directly, with the worst admissible values $A = 2g_K$ and
$\det(K) = 1+A$, over $g_K \le 11$, $p \le 11$ and the full range of $q$ allowed by (F4): 0 failures.

\section{The open case}

Theorem~\ref{thm} does not reach the hyperbolic $L$-space knots, and that is where the inequality is sharp. Baker and
Kegel (arXiv:2203.12013) construct hyperbolic $L$-space knots $K_n$, the closures of the braid words
\[ [(2,1,3,2)^{2n+1},\, -1,\, 2,\, 1,\, 1,\, 2], \]
which are not braid positive for $n=1$ and are not known to be braid positive for $n \ge 2$; Himeno (arXiv:2506.22934) constructs infinitely many hyperbolic $L$-space knots that are provably not braid
positive. Equality $\det = 1 + |\sigma|$ holds for every Baker--Kegel $K_n$ we tested, $n = 1,\dots,25$.

\begin{question}
Does $\det(K) \le 1 + |\sigma(K)|$ hold for every $L$-space knot, in particular for the hyperbolic ones? Equivalently,
by \eqref{eq:det}, is $|\sigma(K)| \ge 2\bigl(\#\{\text{odd gaps}\} - \#\{\text{even gaps}\}\bigr)$?
\end{question}

\paragraph{The sharp case carries a rigidity statement.} If all $g$ gaps are odd then $S \cap [0,2g)$ consists of the
even numbers, so $\Delta_K = \Delta_{T(2,2g+1)}$, $D = g$ and $\det = 2g+1$; Conjecture~\ref{conj} then forces
$|\sigma(K)| = 2g$, that is, the symmetrised Seifert form is definite. So the conjecture implies: every $L$-space knot
with the Alexander polynomial of $T(2,2g+1)$ has $\sigma = -2g$. This is a weak form of Problem 1.21(c)(i) of the K3
problem list \cite{K3} (knot Floer homology detects $T(2g+1,2)$), known there for $g \le 2$ by Ghiggini and by
Farber--Reinoso--Wang. Any proof of Conjecture~\ref{conj} therefore proves this definiteness
statement.

\paragraph{What is known about the sharp case.} Three further facts, contributed by a fourth reader and checked here
against the exact data. (i) Such a $K$ is not a nontrivial cable: the staircase set of a cable $J_{p,q}$ is
$\{ps + qk : s \in S_J,\ 0 \le k \le p-1\}$, whose smallest positive element is $\min(p\,s_{\min}, q) \ge 4$, whereas
$S_{T(2,2g+1)}$ has smallest positive element $2$; here $s_{\min} \ge 2$ because $1$ is a gap of every nontrivial
$L$-space knot, which holds on all 11728 iterated torus knots computed. Among those knots, the only ones with
$\Delta = \Delta_{T(2,2g+1)}$ are the $T(2,2g+1)$ themselves. (ii) The roots of $\Delta_{T(2,2g+1)}$ are simple and on
the unit circle, so the Levine--Tristram function has $g$ jumps of $\pm 2$ on the upper half circle and
$\sigma(K) = -2g + 4m$, where $m$ counts the positive jumps; the question is exactly $m = 0$, that is, that the
signature function of $K$ agrees with that of $T(2,2g+1)$. (iii) Definiteness does not by itself force the torus knot:
among the 11889 $L$-space iterated torus knots computed, exactly 21 are definite, namely the $T(2,n)$ together with
$T(3,4)$ and $T(3,5)$ (the knots among the $ADE$ singularity links). Neither of the last two has the Alexander
polynomial of a $T(2,2g+1)$, so on this restricted class definiteness and detection still coincide.

\paragraph{The hypothesis of the sharp case is met only by $T(2,2g+1)$.} For an $L$-space knot, $\det = 2g+1$ exactly
when every gap is odd, that is when $\Delta = \Delta_{T(2,2g+1)}$, which is also exactly when the knot is thin. Across the
census, the families beyond it, the one-bridge braids, the twisted torus knots and the random search, reaching genus
145 in all, 2009 records meet this hypothesis. Each was put into braid form and found, individually, to be the
closure of $\sigma_1^{2g+1}$ in $B_2$, which identifies it as $T(2,2g+1)$; together they cover $g \le 11$, and
none of the census knots is among them. So every thin $L$-space knot found here is a $T(2,2g+1)$, which is the
$L$-space case of Problem 1.21(c)(i) appearing in the data, and on all of them the sharp case of
Conjecture~\ref{conj} holds.

Together these narrow a counterexample to Conjecture~\ref{conj} in the sharp case: it would have to be hyperbolic or a
satellite with a non-cable pattern, it could not be braid positive, and it would need a positive Levine--Tristram jump
at one of the $g$ roots.

\paragraph{Why the induction stops.} For a torus knot the gaps of $\langle p,q \rangle$ satisfy
$2g - |\sigma| = 4\,\#\{\text{gaps} > pq/2\}$, so Theorem~\ref{thm} for torus knots says
$\#\{\text{gaps} > pq/2\} \le \#\{\text{even gaps}\}$. No analogue of the half-conductor $pq/2$ is available for a
hyperbolic $L$-space knot, and that is exactly what the cabling induction replaces.

\paragraph{Where equality occurs.} Among the hyperbolic $L$-space knots tested, equality occurs exactly on the
Baker--Kegel family: on every $K_n$ tested, and on no other census knot (Himeno's $K_2$ is the same knot as $K_1$). Failing
to be braid positive does not force it: Himeno's $K_3, \dots, K_{8}$ provably are not braid positive and none of them
attains equality, the last having $\det = 39$ against $1 + |\sigma| = 99$. The question is
what distinguishes the Baker--Kegel family.

Two further questions. Which $L$-space knots attain equality? The known ones are $T(2,n)$, the $(2,q)$-cables, the
Baker--Kegel family, and 70 of the 274 one-bridge parameter triples computed here. Does the inequality follow from a property of
the symmetrised Seifert form of an $L$-space knot, its determinant being bounded by its signature? Definite strongly
quasipositive links are studied in \cite{BBG}.

\section{A second statement from the same run}

The program produced one other statement that its checks could not refute and that we could not place in the
literature: for every $L$-space knot, $u(K) = g(K)$, where $u$ is the unknotting number.

Half of it is immediate. For an $L$-space knot $g = g_4$, and $u \ge g_4$ always, so $u \ge g$. The content is
$u \le g$, which holds for braid positive knots by Rudolph, hence for every $L$-space knot known to be braid
positive; by Baker and Kegel \cite{BK} that is all but one of the census $L$-space knots. What is open is the
non-braid-positive case, and there the only evidence is explicit unknottings.

Each certificate is a sequence of $g$ crossing changes ending at the unknot, found by a search that keeps only the
changes lowering $|\tau|$ by exactly one, which is necessary for an unknotting of that length. Certificates exist
for every knot tested where the statement has content: o9\_30634 (Baker--Kegel $K_1$, also Himeno $K_2$), which is not
braid positive; Baker--Kegel $K_2, \dots, K_{10}$, not known to be braid positive, up to genus 42; and
Himeno $K_3$ and $K_4$, which provably are not. That is 12 knots, 3 of them provably outside
Rudolph's theorem, with no failure. Certificates were also produced for 59 census knots and for the $L$-space
twisted torus knots with a non-positive braid word, but the hyperbolic ones among those are braid positive census
knots and the rest are not hyperbolic, so they test the search rather than the statement.

Unlike the inequality, this statement has not been shown to need its hypothesis. The half we can prove, $u \ge g$,
uses the $L$-space property only through $g = g_4$, which already holds for every fibered strongly quasipositive
knot; and we know of no fibered strongly quasipositive knot with $u > g$. Whether $u = g$ is really a statement
about $L$-space knots, or about a wider class, is part of the question.

\section{Reproducibility}\label{sec:repro}

Scripts, each writing the JSON or JSONL file the corresponding numbers are read from:
\texttt{iterated.py} (exact iterated torus knots), \texttt{lemma\_L.py} (all steps of Lemma~\ref{lem:L}),
\texttt{induction\_proof.py} ((F1)--(F3) and the induction step), \texttt{proof\_steps.py} (identity~\eqref{eq:det}),
\texttt{census\_lspace\_fast.sage} (all census $L$-space knots from the Baker--Kegel braid words),
\texttt{census\_sage\_native.sage} (the non-$L$-space comparison), \texttt{families\_fast.sage} (families beyond the
census), \texttt{exact\_signature.py} and \texttt{exact\_census.py} (exact inertia for all census knots),
\texttt{braid\_identity.py} (the closures are the named census knots), \texttt{onebridge.sage} and
\texttt{onebridge\_lspace.py} (the one-bridge braid family), \texttt{twisted\_torus.py} and
\texttt{twisted\_torus\_sigma.sage} (the twisted torus family and its exact signatures), and
\texttt{twisted\_torus\_ug.py} (the unknotting certificates of Section~8), \texttt{definiteness.py} and
\texttt{definiteness\_identify.py} (the sharp case across all families), \texttt{random\_lspace.py} and
\texttt{random\_lspace\_sigma.sage} (the random search), and \texttt{distinct\_knots.py} (how many distinct knots each
generated family contains, and which of its hyperbolic knots are in the census).
Tools: SnapPy 3.3.2 with \texttt{knot\_floer\_homology}, Sage 10.7, khoca 1.5. The Lean formalisation of the
lattice inequality is \texttt{K33Lattice/Basic.lean}, built against Mathlib at revision
\texttt{0df444a360eaa60ab8c11dca51a86af692955474} with \texttt{leanprover/lean4:v4.33.1}.

\paragraph{Data provenance.} The candidate was generated from KnotInfo \cite{KnotInfo}, accessed on 14 September 2026.
No KnotInfo data is reproduced here: its maintainers ask that the database not be reposted, since copies go out
of date, so it should be read from \url{https://knotinfo.org}. The braid words for census knots are from Baker and
Kegel's appendix \cite{BK}.

\begin{thebibliography}{9}
\bibitem{BK} K. L. Baker, M. Kegel, \emph{Census $L$-space knots are braid positive, except for one that is not},
Algebr. Geom. Topol. 24 (2024); arXiv:2203.12013.
\bibitem{BBG} M. Boileau, S. Boyer, C. McA. Gordon, \emph{On definite strongly quasipositive links and $L$-space
branched covers}, Adv. Math. 357 (2019).
\bibitem{EN} D. Eisenbud, W. Neumann, \emph{Three-dimensional link theory and invariants of plane curve
singularities}, Ann. of Math. Studies 110, Princeton Univ. Press (1985).
\bibitem{Feller} P. Feller, \emph{The signature of positive braids is linearly bounded by their genus},
Internat. J. Math. 26 (2015); arXiv:1311.1242.
\bibitem{GL} J. E. Greene, L. Liechti, \emph{On the signature of a positive braid},
Ann. Henri Lebesgue 7 (2024), 823--839; arXiv:2308.02275.
\bibitem{GLM} C. McA. Gordon, R. A. Litherland, K. Murasugi, \emph{Signatures of covering links},
Canad. J. Math. 33 (1981), 381--394.
\bibitem{KnotInfo} C. Livingston and A. H. Moore, \emph{KnotInfo: Table of Knot Invariants}, knotinfo.org,
September 2026 (accessed 14 September 2026).
\bibitem{Hedden} M. Hedden, \emph{On knot Floer homology and cabling II}, Int. Math. Res. Not. (2009).
\bibitem{Hedden2} M. Hedden, \emph{Notions of positivity and the Ozsv\'ath--Szab\'o concordance invariant},
J. Knot Theory Ramifications 19 (2010), 617--629.
\bibitem{Ni} Y. Ni, \emph{Knot Floer homology detects fibred knots}, Invent. Math. 170 (2007), 577--608.
\bibitem{GLV} J. E. Greene, S. Lewallen, F. Vafaee, \emph{$(1,1)$ $L$-space knots},
Compos. Math. 154 (2018), 918--933.
\bibitem{K3} D. Ruberman, R. \.I. Baykur, R. Kirby (eds.), \emph{K3: A New Problem List in Low-Dimensional Topology},
AMS Math. Surveys and Monographs 295 (2026), Problem 1.21.
\bibitem{Himeno} K. Himeno, \emph{Non-braid positive hyperbolic $L$-space knots}, arXiv:2506.22934.
\bibitem{Hom} J. Hom, \emph{A note on cabling and $L$-space surgeries}, Algebr. Geom. Topol. 11 (2011), 219--223.
\bibitem{Lith} R. A. Litherland, \emph{Signatures of iterated torus knots}, in Topology of Low-Dimensional Manifolds,
Lecture Notes in Math. 722 (1979), 71--84.
\bibitem{OS} P. Ozsv\'ath, Z. Szab\'o, \emph{On knot Floer homology and lens space surgeries},
Topology 44 (2005), 1281--1300.
\bibitem{Murasugi} K. Murasugi, \emph{On a certain numerical invariant of link types},
Trans. Amer. Math. Soc. 117 (1965), 387--422.
\end{thebibliography}

\end{document}
